\honest{No Evolutions for Corporations or Nanodevices}{Eliezer Yudkowsky}{2007}

I open with Perry Metzger: ``So long as there are limited resources and multiple competing actors capable of passing on characteristics, you have selection pressure.'' That is qualitative reasoning, I say. How much selection pressure?

My tool is Price's equation: a characteristic changes according to its covariance with relative fitness. The simple form holds only when the characteristic is inherited. A gene that saves its whole species does not spread, since it helps every genotype equally. \nb{This ``Frodo gene'' comes from ``Evolving to Extinction,'' which Yudkowsky posted the day before but which comes after this post in the book.} It is said that Price killed himself over what his equation implied about altruism, ``though he may have had other issues.''

So limited resources and competing replicators, I conclude, ``are not sufficient to give rise to an evolution.'' \nb{By ``an evolution'' the post means cumulative selection that builds complex adaptations. On the standard account, Lewontin's of 1970, variation, differential fitness and heritable fitness are enough for some evolutionary change, so, given variation among the competitors, Metzger's claim that there is selection pressure is true. The question is how much, which is the question the post asked.}

Corporations compete, have offspring and die, but a spinoff resembles its parent only weakly, and CEOs ``cannot divide themselves by fission.'' \nb{Evolutionary economists measure a firm's heredity differently. Nelson and Winter (1982) treat its routines as genes, inherited within the firm over time. The post does not discuss this work.} Failed firms are weeded out, as short-lived stars vanish from view, but one-off weeding builds no complex features. Those come from human design, with bankruptcy adding ``a handful of additional bits.'' I do not calculate the bits.

Nanodevices could copy themselves exactly. ``Correlation is covariance divided by variance,'' I explain, and what drives evolution is covariance, which shrinks when a characteristic barely varies. \nb{Correlation is covariance divided by the product of the two standard deviations. The point about covariance survives the correction.} The Foresight Institute proposes encrypting a nanodevice's instructions so that any copying error scrambles them. Heritable assembly errors, ``prions,'' might remain, but I suppose one that speeds copying from 1000 seconds to 999.99999 and conclude that it would barely spread. Such devices would run out of atoms within a few generations, so grey goo that ate the Earth would then barely change. That, to me, is the most frightening thing about it.

I end with the conditions for cumulative selection: replication, variation in traits and in reproduction, their correlation, faithful long-range heredity, frequent births, and all of it over many iterations. \nb{The first items restate Lewontin's three conditions. The quantitative ones are the post's addition, and they carry its central distinction between one-off and cumulative selection.}
